\appendix

\section{Kalman Filter Design Validation}
\label{sec:kalman_ablation}

This appendix presents ablation studies validating our Kalman filter design choices. All experiments use 21-fold Leave-One-Subject-Out Cross-Validation (LOSO-CV) on the SmartFallMM young subject cohort.

\subsection{Adaptive Measurement Noise}

The accelerometer-derived roll and pitch observations assume that measured acceleration is dominated by gravity. During fall impacts, linear acceleration can spike to 20--30 m/s$^2$, potentially contaminating the gravity reference. We tested adaptive measurement noise gating that inflates $R_{\text{acc}}$ when $\|\mathbf{a}\| > 1.5g$:
\begin{equation}
R_{\text{acc}}^{\text{adaptive}} = \begin{cases}
R_{\text{acc}} \cdot \left(\frac{\|\mathbf{a}\|}{g}\right)^2 & \text{if } \|\mathbf{a}\| > 1.5g \\
R_{\text{acc}} & \text{otherwise}
\end{cases}
\end{equation}

\begin{table}[htbp]
\centering
\caption{Adaptive Measurement Noise Ablation (21-fold LOSO-CV)}
\label{tab:adaptive_racc}
\begin{tabular}{@{}lcc@{}}
\toprule
\textbf{Configuration} & \textbf{Test F1 (\%)} & \textbf{$\Delta$} \\
\midrule
Constant $R_{\text{acc}} = 0.05$ (baseline) & 90.65 & --- \\
Adaptive $R_{\text{acc}}$ (threshold 1.5$g$) & 90.63 & $-0.02$ \\
\bottomrule
\end{tabular}
\end{table}

The negligible difference ($\Delta = -0.02\%$) confirms that constant measurement noise is sufficient for our application. Fall impacts typically span only 3--9 samples (0.1--0.3 seconds) within our 128-sample windows, representing less than 7\% of each window.

\subsection{Yaw Channel Contribution}

Without a magnetometer, yaw ($\psi$) lacks an absolute reference and accumulates drift through gyroscope integration. Typical MEMS gyroscope bias drift is 0.1--1.0$^\circ$/s. Over our 4.27-second windows, maximum drift is bounded to approximately 4$^\circ$, which is negligible compared to fall rotation magnitudes (90--180$^\circ$).

\begin{table}[htbp]
\centering
\caption{Yaw Channel Ablation (21-fold LOSO-CV)}
\label{tab:yaw_ablation}
\begin{tabular}{@{}lcc@{}}
\toprule
\textbf{Configuration} & \textbf{Test F1 (\%)} & \textbf{$\Delta$} \\
\midrule
6-channel (no yaw): [SMV, $\mathbf{a}$, $\phi$, $\theta$] & 89.40 & $-1.25$ \\
7-channel (with yaw): [SMV, $\mathbf{a}$, $\phi$, $\theta$, $\psi$] & 90.65 & --- \\
\bottomrule
\end{tabular}
\end{table}

The 1.25\% F1 improvement with yaw demonstrates that \textit{relative} yaw dynamics within each window provide discriminative information for fall classification, even without absolute heading reference.

\subsection{Gyroscope Bias Estimation}

Standard practice in IMU orientation estimation is to augment the filter state with gyroscope bias terms. We compared our 6-dimensional Linear Kalman Filter against a 9-dimensional Extended Kalman Filter (EKF) with explicit bias estimation:
\begin{equation}
\mathbf{x}_{\text{EKF}} = \begin{bmatrix} \phi & \theta & \psi & \dot{\phi} & \dot{\theta} & \dot{\psi} & b_{\omega_x} & b_{\omega_y} & b_{\omega_z} \end{bmatrix}^T
\end{equation}

\begin{table}[htbp]
\centering
\caption{Filter Complexity Ablation (21-fold LOSO-CV)}
\label{tab:bias_ablation}
\begin{tabular}{@{}lccc@{}}
\toprule
\textbf{Filter} & \textbf{Bias Est.} & \textbf{Test F1 (\%)} & \textbf{$\Delta$} \\
\midrule
Linear Kalman (6D) & No & 90.91 & $+2.02$ \\
Extended Kalman (9D) & Yes & 88.89 & --- \\
\bottomrule
\end{tabular}
\end{table}

Counter-intuitively, the EKF with bias estimation performed \textit{worse} than the simpler Linear Kalman Filter. We attribute this to three factors:
\begin{enumerate}
\item Gyroscope bias is only observable during quasi-static periods, which are rare in 4.27-second windows containing falls
\item The EKF attempts to estimate bias from insufficient data, introducing estimation noise
\item The neural network learns to compensate for consistent filter behavior, whereas unstable EKF estimates add variance the network cannot model
\end{enumerate}

This finding demonstrates that simpler is better for short-window fall detection: the 4.27-second duration bounds bias-induced drift to levels ($<5^\circ$) less harmful than EKF estimation noise under dynamic conditions.

\subsection{Channel-Aware Normalization}
\label{sec:channel_aware_norm}

We tested three normalization strategies for the 7-channel Kalman-fused features:

\begin{table}[htbp]
\centering
\caption{Normalization Strategy Ablation (21-fold LOSO-CV)}
\label{tab:norm_ablation}
\begin{tabular}{@{}lcc@{}}
\toprule
\textbf{Strategy} & \textbf{Test F1 (\%)} & \textbf{$\Delta$} \\
\midrule
No normalization & 89.62 & $-0.45$ \\
Normalize all channels & 89.84 & $-0.23$ \\
Channel-aware (acc only) & 91.10 & --- \\
\bottomrule
\end{tabular}
\end{table}

Channel-aware normalization---applying StandardScaler only to accelerometer channels while preserving orientation angles in raw radians---achieves the best performance. This design reflects the different characteristics of each modality:
\begin{itemize}
\item Accelerometer channels can spike to 30+ m/s$^2$ during falls, requiring normalization to prevent scale domination
\item Orientation channels are already bounded ($\pm\pi$ radians) with preserved physical meaning
\end{itemize}

\section{Architecture Ablation Studies}
\label{sec:arch_ablation}

\subsection{Single-Stream vs Dual-Stream}

We compared our dual-stream architecture against a single-stream baseline processing all 7 channels through a shared embedding:

\begin{table}[htbp]
\centering
\caption{Stream Architecture Ablation (21-fold LOSO-CV)}
\label{tab:stream_ablation}
\begin{tabular}{@{}lcc@{}}
\toprule
\textbf{Architecture} & \textbf{Test F1 (\%)} & \textbf{Val-Test Gap} \\
\midrule
Single-stream (7ch shared) & 89.49 & 6.61\% \\
Dual-stream (4ch + 3ch) & 91.10 & 5.00\% \\
\bottomrule
\end{tabular}
\end{table}

The dual-stream architecture improves both test F1 (+1.61\%) and generalization (smaller val-test gap). Separating accelerometer and orientation streams prevents noisy orientation estimates from corrupting acceleration feature extraction.

\subsection{Embedding Dimension Ratio}

The dual-stream architecture allocates embedding dimensions between accelerometer and orientation streams. We tested different allocation ratios:

\begin{table}[htbp]
\centering
\caption{Embedding Ratio Ablation (21-fold LOSO-CV)}
\label{tab:ratio_ablation}
\begin{tabular}{@{}lcc@{}}
\toprule
\textbf{Ratio (Acc:Ori)} & \textbf{Test F1 (\%)} & \textbf{$\Delta$} \\
\midrule
50:50 (32:32) & 90.23 & $-0.87$ \\
65:35 (42:22) & 91.10 & --- \\
75:25 (48:16) & 90.45 & $-0.65$ \\
\bottomrule
\end{tabular}
\end{table}

The 65:35 ratio allocates more capacity to the accelerometer stream, reflecting its higher signal quality and dimensionality (4 vs 3 channels).

\subsection{SE Module and TAP Contribution}

\begin{table}[htbp]
\centering
\caption{Attention Module Ablation (21-fold LOSO-CV)}
\label{tab:attention_ablation}
\begin{tabular}{@{}lcc@{}}
\toprule
\textbf{Configuration} & \textbf{Test F1 (\%)} & \textbf{$\Delta$} \\
\midrule
No SE, No TAP (mean pool) & 88.92 & $-2.18$ \\
SE only & 89.84 & $-1.26$ \\
TAP only & 90.18 & $-0.92$ \\
SE + TAP (full model) & 91.10 & --- \\
\bottomrule
\end{tabular}
\end{table}

Both SE (channel attention) and TAP (temporal attention) contribute to performance. TAP provides larger gains, suggesting that temporal weighting to focus on fall-critical frames is particularly important.
